\section{Matrizes}

Matrizes são objetos tabulares, dispondo números por linhas e colunas que são referenciados
pela sua posição.
O número de linhas e colunas da matriz definem a sua dimensão, ou tipo.
Essa característica irá restringir quais tipos de operações poderão ser realizadas
envolvendo matrizes.

\begin{def:matriz}
Dados inteiros positivos $m$ e $n$, uma matriz $A$ de dimensão -- ou tipo -- $m \times n$,
é uma tabela com $m$ linhas e $n$ colunas cujos elementos, denominados entradas ou
coeficientes, são números reais; de forma sintética, $A \in \real^{m \times n}$.
A entrada de $A$ da linha de índice $i \in [m]$ e coluna de índice $j \in [n]$ é denotada
como $a_{ij} \in \real$.
\end{def:matriz}

\begin{exp:repr matriz}
Exemplos de matrizes:
\[
	\begin{bmatrix}
		1 & 0 \\
		3 & 2
	\end{bmatrix}, \quad
	\begin{bmatrix}
		5 & 0  & 5  \\
		8 & 2  & 10 \\
		5 & 11 & 4  \\
	\end{bmatrix} \quad \text{e} \quad
	\begin{bmatrix}
		5 & 4  & 3  \\
		0 & 12 & 20 \\
	\end{bmatrix} \quad
\]
\end{exp:repr matriz}

Para que duas matrizes sejam iguais, elas devem ter as mesmas entradas nas mesmas posições,
precisando, pois, serem do mesmo tipo.

\begin{def:igualdade matrizes}
Duas matrizes $A$ e $B$ são iguais se, e somente se, são do mesmo tipo $m \times n$ e
para quaisquer índices $i \in [m]$ e $j \in [n]$, a entrada $a_{ij}$ de $A$ é igual à
entrada $b_{ij}$ de $B$.
\[
	\forall i \in [m]\,\forall j \in [n]\,, ( a_{ij} = b_{ij} )
\]
\end{def:igualdade matrizes}

Um tipo especial de matrizes são aquelas cujo número de linhas é igual a de colunas,
que são chamadas de \textbf{quadradas}.
Em outro caso, elas são chamadas de \textbf{retangulares}, havendo dois subtipos:
as \textbf{altas}, que possuem mais linhas do que colunas, e as \textbf{largas}, as quais
o número de colunas supera o de linhas.

\begin{def:matriz quadrada}
Uma matriz $A$ do tipo $m \times n$ é dita quadrada se $m = n$.
Se $A$ é do tipo $n \times n$, então dizemos que $A$ é de ordem $n$.
\end{def:matriz quadrada}

Um elemento importante das matrizes quadradas são as suas diagonais: a principal, que
consiste daquela que vai do canto suprior esquerdo até o inferior direto, e a secudária,
que vai do elemento suprior direito ao inferior esquerdo.

\begin{def:diagonal}
Seja $A$ uma matriz quadrada $n \times n$.
A digonal principal de $A$ é composta pelas entradas $a_{ij}$ com $i = j$.
A diagonal secudária é composta pelas entradas $a_{ij}$ tal que $i + j = n + 1$.
\end{def:diagonal}

\begin{def:transposta}
Se $A$ é uma matriz $m \times n$, então a sua transposta, denotada por $A^\top$, é
a matriz $n \times m$ cujas entradas $p_{ij}$ são dadas por
\[
	p_{ij} = a_{ji}
\]
\end{def:transposta}

\begin{exp:transposta1}
\[
	A = \begin{bmatrix}
		3 & 2 \\
		1 & 0
	\end{bmatrix} \rightarrow
	A^\top = \begin{bmatrix}
		3 & 1 \\
		2 & 0
	\end{bmatrix}
\]
\end{exp:transposta1}

\begin{exp:transposta2}
\[
	A = \begin{bmatrix}
		2 & 4  & 4  \\
		2 & -2 & 10
	\end{bmatrix} \rightarrow
	A^\top = \begin{bmatrix}
		2 & 2  \\
		4 & -2 \\
		4 & 10
	\end{bmatrix}
\]
\end{exp:transposta2}

Uma classe especial de matrizes são as chamadas matrizes \textbf{esparsas}, que surgem
em diversos problemas aplicados e possuem interessantes propriedades computacionais.
Essas matrizes caracterizam-se por ter boa parte de suas entradas iguais a $0$, porém
não há uma condição que diga de maneira inequívoca se dada matriz seja esparsa ou não.
No entanto, há outras classificações de matrizes que satisfazem determinadas definições que
são amplamente consideradas como esparsa.
Um exemplo trivial de matriz esparsa é matriz nula: aquela cujas entradas são todas zero.

\begin{def:matriz nula}
A matriz $O \in \real^{m \times n}$ é nula se, e somente se, para quaisquer $i \in [m]$
e $j \in [n]$, tem-se que $o_{ij} = 0$.
\end{def:matriz nula}

Outra classificação, de grande importância na Álgebra Linear, são as matrizes diagonais,
cujos elementos fora da diagonal principal são iguais a zero.

\begin{def:matriz diagonal}
Uma matriz quadrada $\Lambda \in \real^{n \times n}$ é dita diagonal se para quaisquer
$i, j \in [n]$ é satisfeito que
\[
	i \ne j \rightarrow \lambda_{ij} = 0
\]
\end{def:matriz diagonal}

\begin{exp:diagonal1}
As matrizes
\[
	\begin{bmatrix}
		2 & 0 \\
		0 & 1
	\end{bmatrix} \quad \text{ e } \quad
	\begin{bmatrix}
		4 & 0  & 0 \\
		0 & -1 & 0 \\
		0 & 0  & 3
	\end{bmatrix}
\]
são diagonais.
\end{exp:diagonal1}

Um exemplo especial de matriz diagonal e, portanto, esparsa, é a matriz \textbf{identidade},
cuja diagonal é preenchida apenas por $1$.

\begin{def:matriz identidade}
Chama-se matriz identidade de ordem $n$, denotada por $\mathbf I_n$, toda matriz diagonal
$n \times n$ cujas entradas da diagonal principal são todas iguais a $1$.
\end{def:matriz identidade}

Em seguida, definimos as matrizes básicas como aquelas cujas entradas são todas iguais
iguais a zero, a menos de uma.

\begin{def:matriz básica}
Uma matriz básica $B_{rs}(\beta)$ é uma matriz quadrada cujas entradas são definidas tal como
\[
	b_{ij} = \begin{cases}
		\beta, & \text{se } i = r \text{ e } j = s \\
		0,     & \text{caso contrário.}
	\end{cases}
\]
\end{def:matriz básica}

\begin{exp:matriz basica}
A matriz $3 \times 3$ a seguir
\[
	B_{21}(1) = \begin{bmatrix}
		0 & 0 & 0 \\
		1 & 0 & 0 \\
		0 & 0 & 0
	\end{bmatrix}
\]
é básica.
\end{exp:matriz basica}

Por sua vez, as matrizes elementares são aquelas cuja a diagonal princpal é toda igual a
$1$ e, caso não seja diagonal, então exatamente um elemento fora da diagonal principal
é não nulo.

\begin{def:matriz elementar}
Uma matriz quadrada $L_{rs}(\alpha)$ é dita elementar se para quaisquer índices $i, j \in [n]$
temos que
\[
	\ell_{ij} = \begin{cases}
		1,      & \text{se } i = j                  \\
		\alpha, & \text{se } i = r \text{ e } j = s \\
		0,      & \text{caso contrário}
	\end{cases}
\]
\end{def:matriz elementar}

Outras matrizes esparsas de destaque na Álgebra Linear são as triangulares, cujo nome
deriva do seu formato tipo escada.
Uma triangular superior é aquela cujas entradas abaixo da diagonal principal são nulos,
enquanto que as inferiores possuem as entradas acima da diagonal principal iguais as zero.

\begin{def:matriz triangular}
Uma matriz quadrada $U$ é dita triangular superior se para quaisquer $i, j \in [n]$
\[
	i > j \rightarrow u_{ij} = 0
\]
Uma matriz quadrada $L$ é traingular inferior se para quaisquer $i, j \in [n]$
\[
	i < j \rightarrow u_{ij} = 0
\]
\end{def:matriz triangular}

\begin{exp:triangular1}
A matriz
\[
	\begin{bmatrix}
		2 & 3 & 0 \\
		0 & 4 & 1 \\
		0 & 0 & 8
	\end{bmatrix}
\]
é triangular superior
\end{exp:triangular1}

\begin{exp:triangular2}
A matriz
\[
	\begin{bmatrix}
		0  & 0 & 0 \\
		2  & 4 & 0 \\
		-5 & 1 & 8
	\end{bmatrix}
\]
é triangular inferior.
\end{exp:triangular2}

\begin{def:simetrica}
Uma matriz $A$ é dita simétrica se, e somente se, para qualquer índices $i$ e $j$,
\[
	a_{ij} = a_{ji}
\]
o que equivale a
\[
	A^\top = A
\]
\end{def:simetrica}

\begin{exp:simetrica1}
As matrizes identidade e nula são matrizes simétricas.
\end{exp:simetrica1}

\begin{exp:simetrica2}
A matriz
\[
	\begin{bmatrix}
		2  & -3 & 7 \\
		-3 & 1  & 9 \\
		7  & 9  & 5
	\end{bmatrix}
\]
é simétrica.
\end{exp:simetrica2}
